\documentclass[10pt,a4paper]{amsart}
\usepackage[margin=19mm]{geometry}
\usepackage{amsmath,amssymb,mathtools}
\usepackage[scr=rsfs,cal=euler]{mathalfa}
\usepackage{xfrac}
\usepackage[unicode,hidelinks,bookmarks=false]{hyperref}
\newcommand{\ie}{i.e.\ }
\newcommand{\cf}{cf.\ }
\newcommand{\resp}{respectively\ }
\newcommand{\Subsection}{\subsection}
\providecommand{\nobreakdash}{\nobreak\hbox{-}}
\setlength{\emergencystretch}{2em}
\multlinegap0pt
\usepackage{bm}
\usepackage{amssymb}
\usepackage[shortlabels]{enumitem}
\usepackage{xcolor}
\usepackage{tikz}
\newtheorem{assumption}{Assumption}
\newtheorem{theorem}{Theorem}
\DeclareMathOperator*{\argmin}{arg\,min}
\def\cC{{\mathcal{C}}}
\def\cF{{\mathcal{F}}}
\def\eqref#1{equation~\ref{#1}}
\def\mA{{\bm{A}}}
\def\mB{{\bm{B}}}
\def\mI{{\bm{I}}}
\def\mS{{\bm{S}}}
\def\sM{{\mathbb{M}}}
\def\sN{{\mathbb{N}}}
\def\sR{{\mathbb{R}}}
\def\sT{{\mathbb{T}}}
\def\sU{{\mathbb{U}}}
\def\sX{{\mathbb{X}}}
\def\vh{{\bm{h}}}
\def\vn{{\bm{n}}}
\def\vp{{\bm{p}}}
\def\vu{{\bm{u}}}
\def\vv{{\bm{v}}}
\def\vx{{\bm{x}}}
\def\vy{{\bm{y}}}
\def\vz{{\bm{z}}}
\def\vzero{{\bm{0}}}
\title{Expressive Power of Implicit Models: Rich Equilibria and Test-Time Scaling}
\author{Jialin Liu, Lisang Ding, Stanley Osher, Wotao Yin}
\date{Source: arXiv:2510.03638v4 (CC BY 4.0)}
\begin{document}
\maketitle

\noindent\emph{Source and licence.} Expressive Power of Implicit Models: Rich Equilibria and Test-Time Scaling. Version arXiv:2510.03638v4 (CC BY 4.0), distributed by arXiv under the Creative Commons Attribution 4.0 International licence, http://creativecommons.org/licenses/by/4.0/, as recorded in the arXiv licence metadata for this exact version and shown on the abstract page https://arxiv.org/abs/2510.03638v4. Full licence notice: \texttt{licenses/arxiv-2510.03638-CC-BY-4.0.txt}.\par\smallskip

\noindent\textit{Editorial scope.} Counted unit: the article's theorem thm:inverse-lip-1b (source0.tex lines 863--866). The article's complete original proof of it is delivered with it (source0.tex lines 1980--2137, one \texttt{proof} environment of 158 lines, which the article places away from the statement). No mathematics is written, repaired or re-derived: the statement and the proof are the article's own lines, in the article's own notation, and no retained line is substituted or re-flowed.  Retained uncounted as the source-local context the proof needs: lines 831--834; lines 836--839; lines 844--847; lines 858--861.  Nothing was omitted from any retained span, so the package carries no omission record.  No retained line contains a citation, so the wrapper carries no bibliography and no bibliography entry.  No retained line contains a figure environment, an image-inclusion command or drawing code, so no image is delivered and the package declares no asset dependency.  Editorial formatting only, in addition: an AMS \texttt{amsart} preamble that restates the article's own declarations and macros used by the retained lines, namely the environments \texttt{assumption}, \texttt{theorem} on separate counters, together with the standard packages those lines need.  The retained environments print in this document as Assumption 1, Theorem 1 in order of appearance, whereas the article prints the same items with the numbering of its own full document.  The complete native source of the article as distributed in the arXiv e-print for this exact version is bundled unchanged as \texttt{sources/186375/source0.tex}.
\begin{equation}
\label{eq:mapping-1a}
\min_{\vy\in\sR^n}\;\frac{1}{2}\|\vx-\mA\vy\|^2 \;+\; \frac{\alpha}{2}\,\mathrm{dist}^2(\vy,\sM),
\end{equation}

\begin{equation}
\label{eq:mapping-1b}
\min_{\vy,\vz\in\sR^n}\;\frac{1}{2}\|\vx-\mA\vy\|^2 \;+\; \frac{\alpha}{2}\,\mathrm{dist}^2(\vz,\sM) \;+\; \frac{\beta}{2}\|\vy-\vz\|^2.
\end{equation}

\begin{assumption}
\label{assume:M}
Let $\sM\subset\sR^n$ be a compact, $\cC^2$, embedded (possibly nonconvex) submanifold with positive reach $\tau>0$. Assume the forward operator $\mA:\sR^n\!\to\!\sR^d$ is $(\mu,L)$–bi-Lipschitz when restricted to $\sM$ and let $\sigma_{\textup{max}}$ denote the maximal singular value of $\mA$.
\end{assumption}

\begin{theorem}
\label{thm:inverse-lip-1a}
Under Assumption \ref{assume:M}, there exists $\alpha > 0$ for all $\vx \in \sX$ such that the minimization problem \eqref{eq:mapping-1a} yields a unique minimizer $\hat{\vy}$. Let $\cF_{1a}$: $\vx \mapsto \hat{\vy}$ denote the associated solution map from input $\vx$ to the recovery $\hat{\vy}$. Then $\cF_{1a}$ is locally Lipschitz continuous on $\sX$.
\end{theorem}

\begin{theorem}
\label{thm:inverse-lip-1b}
Under Assumption \ref{assume:M}, there exist $\alpha,\beta > 0$ for all $\vx \in \sX$ such that the minimization problem \eqref{eq:mapping-1b} yields a unique minimizer $(\hat{\vy},\hat{\vz})$. Let $\cF_{1b}$: $\vx \mapsto \hat{\vy}$ denote the associated solution map from input $\vx$ to the recovery $\hat{\vy}$. Then $\cF_{1b}$ is locally Lipschitz continuous on $\sX$.
\end{theorem}

\begin{proof}[Proof of Theorem \ref{thm:inverse-lip-1b}] 
For simplicity, we denote the objective function in \eqref{eq:mapping-1b} as $F_{1b}(\vy,\vz)$:
\[ 
F_{1b}(\vy,\vz) := \frac{1}{2}\|\vx - \mA \vy\|^2 + \frac{\alpha}{2} \textup{dist}^2(\vz,\sM) + \frac{\beta}{2} \|\vz-\vy\|^2,
\]
and we will study its properties analogously to $F_{1a}$.

\textbf{Step 1: Existence of minimizers of $F_{1b}$.} For any $r>0$, as 
\[ \alpha \geq \frac{\|\vn\|^2}{r^2}, \quad \beta \geq \frac{\|\vn\|^2}{r^2}, \]
it holds that
\begin{equation}
\label{eq:proof-thm-mapping-1b-step1}
\inf_{\vy,\vz}F_{1b}(\vy,\vz) = \inf_{ (\vy,\vz): ~ \textup{dist}(\vz,\sM)\leq r \textup{ and } \|\vz-\vy\|\leq r }F_{1b}(\vy,\vz).
\end{equation}
This can be proved by contradiction: (I) Suppose $F_{1b}(\hat{\vy},\hat{\vz})$ is lower than the right-hand-side of \eqref{eq:proof-thm-mapping-1b-step1} and $\textup{dist}(\hat\vz,\sM) > r$, we have
\[ F_{1b}(\hat{\vy},\hat{\vz}) \geq 0 + \frac{\|\vn\|^2}{2 r^2} \textup{dist}^2(\hat{\vz},\sM) + 0 > \frac{1}{2} \|\vn\|^2 = F_{1b}(\vy_\ast,\vy_\ast)\]
which contradicts with the hypothesis regarding $(\hat{\vy},\hat{\vz})$. (II)  Suppose $F_{1b}(\hat{\vy},\hat{\vz})$ is lower than the right-hand-side of \eqref{eq:proof-thm-mapping-1b-step1} and $\|\hat{\vz}-\hat{\vy}\| > r$, we have
\[ F_{1b}(\hat{\vy},\hat{\vz}) \geq 0 + 0 + \frac{\|\vn\|^2}{2 r^2}\|\hat{\vz}-\hat{\vy}\|^2  > \frac{1}{2}\|\vn\|^2 = F_{1b}(\vy_\ast,\vy_\ast)\]
which also derives a contradiction. Arguments in (I) and (II) together prove \eqref{eq:proof-thm-mapping-1b-step1}. Similar to the proof of Theorem \ref{thm:inverse-lip-1a}, \eqref{eq:proof-thm-mapping-1b-step1} implies the existence of minimizers of $F_{1b}$ (i.e., minimizers are attainable.)

\textbf{Step 2: Bound of minimizers of $F_{1b}$.} To extend the proof regarding $F_{1a}$ to $F_{1b}$, we consider the following inequality that holds for all $\vy,\vz \in \sU_{\tau}(\sM)$
\[ \|\vy-\vp(\vy)\| \leq \|\vy-\vp(\vz)\| \leq \|\vy-\vz\| + \|\vz-\vp(\vz)\| = \|\vy-\vz\| + \textup{dist}(\vz,\sM) \leq 2r. \]
Therefore, we need $2r < \tau$ and
\begin{equation}
\label{eq:proof-alpha-3}
  \alpha \geq \frac{\|\vn\|^2}{r^2} > \frac{4\|\vn\|^2}{\tau^2}, \quad \beta \geq \frac{\|\vn\|^2}{r^2}  > \frac{4\|\vn\|^2}{\tau^2} 
\end{equation}
to ensure $\hat{\vy},\hat{\vz} \in \sU_{\tau}(\sM)$. Following the same argument as the proof of Theorem \ref{thm:inverse-lip-1a}, the above condition \eqref{eq:proof-alpha-3} implies
\[
\|\vp(\hat{\vy}) - \vy_\ast \| \leq \frac{\sigma_{\textup{max}}}{\mu} (2r) + \frac{2}{\mu} \|\vn\| 
\]
and hence
\begin{equation}
\label{eq:proof-error-bound-1b}
\|\hat{\vy} - \vy_\ast\| \leq \|\hat{\vy} - \vp(\hat{\vy}) \| + \|\vp(\hat{\vy}) - \vy_\ast \| \leq 2 \left(1 + \frac{\sigma_{\textup{max}}}{\mu} \right) r + \frac{2}{\mu} \|\vn\|
\end{equation}
and
\begin{equation}
\label{eq:proof-error-bound-1b-z}
\|\hat{\vz} - \vy_\ast\| \leq \|\hat{\vz} - \hat{\vy} \| + \|\hat{\vy} - \vy_\ast \| \leq \left(3 + 2 \frac{\sigma_{\textup{max}}}{\mu} \right) r + \frac{2}{\mu} \|\vn\|
\end{equation}
holds for all $(\hat{\vy},\hat{\vz})$ that minimizes $F_{1b}(\vy,\vz)$.

\textbf{Step 3: Positive definiteness of the Hessian of $F_{1b}$.} Function $F_{1b}(\vy,\vz)$'s Hessian matrix is of size $2n \times 2n$ and can be written as a $2\times 2$ block w.r.t. $\vy$ and $\vz$:
\[ \nabla^2F_{1b}(\vy,\vz) = 
\begin{bmatrix}
\mA^\top \mA & \vzero \\ 
\vzero & \vzero
\end{bmatrix} + 
\alpha \begin{bmatrix}
\vzero & \vzero \\ 
\vzero & \mI - D\vp(\vz)
\end{bmatrix} + 
\beta \begin{bmatrix}
\mI & -\mI \\ 
-\mI & \mI
\end{bmatrix}
\]
For any $\vh = [\vu^\top~~ \vv^\top]^\top \in \sR^{2n}$, the quadratic form $\langle \vh, \nabla^2F_{1b}(\vy,\vz) \vh\rangle$ can be calculated through:
\[
\langle \vh, \nabla^2F_{1b}(\vy,\vz) \vh\rangle
= \vu^\top\mA^\top \mA\vu + \alpha \vv^\top(\mI - D\vp(\vz))\vv + \beta \|\vu-\vv\|^2
\]
Decompose $\vu = \vu_{\mathrm{T}} + \vu_{\mathrm{N}}$ and $\vv = \vv_{\mathrm{T}} + \vv_{\mathrm{N}}$ in $\sT_{\vp(\vz)}(\sM) \oplus \sN_{\vp(\vz)}(\sM)$. Using the same argument as the proof of Theorem \ref{thm:inverse-lip-1a}, we have
\begin{align*}
& \langle \vh, \nabla^2F_{1b}(\vy,\vz) \vh\rangle \\
\geq & \bigg( \mu^2 \|\vu_{\mathrm{T}} \|^2 - 2 \sigma_{\textup{max}} L\|\vu_{\mathrm{T}}\|\|\vu_{\mathrm{N}}\| \bigg) 
+ \alpha \left(- \frac{s}{\tau - s} \|\vv_{\mathrm{T}} \|^2 + \|\vv_{\mathrm{N}} \|^2  \right)
+ \beta \|\vu-\vv\|^2
\end{align*}
which implies 
\begin{align*}
& \langle \vh, \nabla^2F_{1b}(\vy,\vz) \vh\rangle  \\
\geq & \bigg( \mu^2 \|\vu_{\mathrm{T}} \|^2 - 2 \sigma_{\textup{max}} L\|\vu_{\mathrm{T}}\|\|\vu_{\mathrm{N}}\| \bigg) \\
& \quad + \alpha \left(- \frac{s}{\tau - s} \|\vv_{\mathrm{T}} \|^2 + \|\vv_{\mathrm{N}} \|^2  \right)
+ \beta \bigg( \|\vu_{\mathrm{T}}-\vv_{\mathrm{T}}\|^2 + \|\vu_{\mathrm{N}}-\vv_{\mathrm{N}}\|^2 \bigg) \\
\geq & \bigg( \mu^2 \|\vu_{\mathrm{T}} \|^2 - 2 \sigma_{\textup{max}} L\|\vu_{\mathrm{T}}\|\|\vu_{\mathrm{N}}\| \bigg) \\
& \quad + \alpha \left(- \frac{s}{\tau - s} \|\vv_{\mathrm{T}} \|^2 + \|\vv_{\mathrm{N}} \|^2  \right)
+ \beta \bigg( ( \|\vu_{\mathrm{T}}\| - \|\vv_{\mathrm{T}}\| )^2 + ( \|\vu_{\mathrm{N}}\| - \| \vv_{\mathrm{N}}\| )^2 \bigg) \\
= & \begin{bmatrix}
\|\vu_{\mathrm{T}}\| & \|\vu_{\mathrm{N}}\| & \|\vv_{\mathrm{T}}\| & \|\vv_{\mathrm{N}}\|
\end{bmatrix}
\underbrace{\begin{bmatrix}
\mu^2 + \beta & - \sigma_{\textup{max}} L & - \beta & \\
- \sigma_{\textup{max}} L & \beta & & -\beta \\
-\beta & & \beta - \alpha \frac{s}{\tau-s} & \\
& - \beta & & \alpha + \beta 
\end{bmatrix}}_{=:\mB}
\begin{bmatrix}
\|\vu_{\mathrm{T}}\| \\ \|\vu_{\mathrm{N}}\| \\ \|\vv_{\mathrm{T}}\| \\ \|\vv_{\mathrm{N}}\|
\end{bmatrix}
\end{align*}
To ensure the positive definiteness of $\nabla^2F_{1b}(\vy,\vz)$, it's enough to ensure $\mB \succ \vzero$. For simplicity, we define
\[
\theta := \alpha \frac{s}{\tau-s}, \quad 
\mB_{1}:= \begin{bmatrix}
\mu^2 + \beta & - \sigma_{\textup{max}} L \\
- \sigma_{\textup{max}} L & \beta 
\end{bmatrix}
\quad \mB_{2}:= \begin{bmatrix}
 - \beta & \\
 & -\beta 
\end{bmatrix}
\quad \mB_{3}:= \begin{bmatrix}
 \beta - \theta & \\
 & \alpha + \beta 
\end{bmatrix}\]
Then $\mB = \begin{bmatrix}
\mB_1 & \mB_2 \\ 
\mB_2^\top & \mB_3
\end{bmatrix}$ is positive definite if and only if $\mB_3$ and its Schur complement $\mS$ are both positive definite:
\[ \mB_3 \succ \vzero ,\quad \mS = \mB_1 - \mB_2 \mB_3^{-1} \mB_2^\top \succ \vzero\] 
As $\mB_2$ and $\mB_3$ are both diagonal, so $\mB_2 \mB_3^{-1} \mB_2^\top$ is straight forward to calculate: $\mB_2 \mB_3^{-1} \mB_2^\top = \textup{diag}\left(\frac{\beta^2}{\beta-\theta}, \frac{\beta^2}{\alpha+\beta}\right)$. Then the Schur complement can be calculated:
\[ \mS =  \begin{bmatrix}
\mu^2 + \beta - \frac{\beta^2}{\beta-\theta} & - \sigma_{\textup{max}} L \\
- \sigma_{\textup{max}} L & \beta - \frac{\beta^2}{\alpha+\beta}
\end{bmatrix} = 
\begin{bmatrix}
\mu^2 - \frac{\beta \theta}{\beta-\theta} & - \sigma_{\textup{max}} L \\
- \sigma_{\textup{max}} L & \frac{\alpha\beta}{\alpha+\beta}
\end{bmatrix}
\]
Note that $\mB_3 \succ \vzero$ if.f $\beta > \theta$. Therefore, $\mB \succ \vzero$ if.f.
\begin{equation}
\label{eq:proof-B-PD}
\beta > \theta, \quad \mu^2 > \frac{\beta \theta}{\beta-\theta}, \quad \left(\mu^2 - \frac{\beta \theta}{\beta-\theta}\right)\frac{\alpha\beta}{\alpha+\beta} > \sigma^2_{\textup{max}} L^2,
\end{equation}
where $\theta = \alpha \frac{s}{\tau-s}$. Finally, we obtain that \eqref{eq:proof-B-PD} ensures $\nabla^2F_{1b}(\vy,\vz)\succ \vzero$ for all $\vy \in \sR^n$ and all $\vz \in \overline{\sU}_s(\sM)$ with $s < \tau$.

\textbf{Step 4: Uniqueness of minimizers of $F_{1b}$.} Comparable to the Step 4 in Theorem \ref{thm:inverse-lip-1a}, we need $\|\hat{\vz}-\vy_\ast\| \leq s$ for all $(\hat{\vy},\hat{\vz}) \in \argmin F_{1b}(\vy,\vz)$. Based on \eqref{eq:proof-error-bound-1b-z}, it's enough to guarantee
\begin{equation}
\label{eq:proof-1b-s-condition}
\left(3 + 2 \frac{\sigma_{\textup{max}}}{\mu} \right) r + \frac{2}{\mu} \|\vn\| \leq s
\end{equation}
Now we choose
\[ s = \frac{4}{\mu} \|\vn\|, \quad r = \frac{2}{5\sigma_{\textup{max}}} \|\vn\|\]
which directly satisfies \eqref{eq:proof-1b-s-condition}. As $\|\vn\| < \frac{1}{76}\frac{\mu^5}{\sigma^2_{\textup{max}}L^2} \tau$, we have
\[ s = \frac{4}{\mu} \|\vn\| < \frac{1}{19} \frac{\mu^4}{\sigma^2_{\textup{max}}L^2} \tau, \quad \frac{s}{\tau-s} < \frac{\frac{1}{19} \frac{\mu^4}{\sigma^2_{\textup{max}}L^2} \tau}{\tau - \frac{1}{19} \frac{\mu^4}{\sigma^2_{\textup{max}}L^2} \tau} \leq \frac{\frac{1}{19} \frac{\mu^4}{\sigma^2_{\textup{max}}L^2} \tau}{\tau - \frac{1}{19}  \tau} = \frac{1}{18} \frac{\mu^4}{\sigma^2_{\textup{max}}L^2}\]
As long as we take
\[\alpha = \frac{9\sigma^2_{\textup{max}}L^2}{\mu^2}, \quad \beta \geq \max\left(\alpha, \frac{3}{2} \mu^2\right)\]
it holds that
\[
\theta = \alpha \frac{s}{\tau-s} < \frac{9\sigma^2_{\textup{max}}L^2}{\mu^2} \frac{1}{18} \frac{\mu^4}{\sigma^2_{\textup{max}}L^2} = \frac{1}{2} \mu^2 
\]
which implies $\beta > 3 \theta$ and hence $\beta > \theta$. Moreover, we can verify the remaining part of \eqref{eq:proof-B-PD}:
\[
\begin{aligned}
\frac{\beta \theta}{\beta-\theta} <&  \frac{\beta \theta}{\beta - \beta/3} = \frac{3}{2}\theta < \frac{3}{4}\mu^2 < \mu^2, \\ 
\left(\mu^2 - \frac{\beta \theta}{\beta-\theta}\right)\frac{\alpha\beta}{\alpha+\beta} >&  \left(\mu^2 - \frac{3}{4}\mu^2\right) \frac{\alpha\beta}{\beta+\beta} = \frac{1}{8} \mu^2 \alpha = \frac{1}{8} \mu^2 \cdot \frac{9\sigma^2_{\textup{max}}L^2}{\mu^2} > \sigma^2_{\textup{max}}L^2.
\end{aligned}
\]
which finishes the proof of \eqref{eq:proof-B-PD}. Finally, it's enough to verify \eqref{eq:proof-alpha-3}:
  \[2r \leq \frac{\|\vn\|}{\sigma_{\textup{max}}} \leq \frac{1}{76}\frac{\mu^5}{\sigma^3_{\textup{max}}L^2} \tau < \tau, \quad \frac{\|\vn\|^2}{r^2} = \frac{25}{4} \sigma^2_{\textup{max}} \leq \alpha \leq \beta,\]
  which finishes Step 4, and concludes the uniqueness of $(\hat{\vy},\hat{\vz})$.
  

\textbf{Step 5: Local Lipschitz continuity of $\cF_{1b}$.} By largely following Step 5 in the proof of Theorem \ref{thm:inverse-lip-1a} and changing $\nabla^2 F_{1a}(\vy)$ to $\nabla^2 F_{1b}(\vy,\vz)$, one can directly conclude that the mapping $\cF_{1b}$ is locally Lipschitz continuous on $\sX$.
\end{proof}

\end{document}
